\startSection
	\selectlanguage{spanish}
	\sectionTitle{Himno Nacional de Bolivia}
	
	\begin{question}{En el contexto de la historia musical boliviana, ?`quién fue el compositor responsable de la música del Himno Nacional de Bolivia, pieza que ha acompañado actos oficiales desde el siglo XIX?}
		\choice{José Ignacio de Sanjinés}
		\choice{Eduardo Abaroa}
		\choice[correct]{Leopoldo Benedetto Vincenti}
		\choice{Franz Tamayo}
		\choice{Modesto Omiste}
	\end{question}

	
	\begin{question}{La letra del Himno Nacional de Bolivia fue escrita por un destacado intelectual boliviano, reflejando el espíritu patriótico de la época. ?`Quién fue el autor de esta letra?}
		\choice{Ricardo José Bustamante}
		\choice{Franz Tamayo}
		\choice{Adela Zamudio}
		\choice{Manuel José Cortés}
		\choice[correct]{José Ignacio de Sanjinés}
	\end{question}

	
	\begin{question}{El Himno Nacional de Bolivia fue oficialmente adoptado en una fecha significativa para la consolidación de los símbolos patrios. ?`En qué año se oficializó su uso como himno nacional?}
		\choice{1825}
		\choice{1842}
		\choice[correct]{1845}
		\choice{1851}
		\choice{1879}
	\end{question}

	
	\begin{question}{El Himno Nacional de Bolivia comienza con un verso que evoca el destino favorable de la nación. ?`Cuál es el primer verso que da inicio a esta emblemática composición?}
		\choice{\enquote{De la Patria, el alto nombre}}
		\choice{\enquote{Compatriotas, la Patria muriendo}}
		\choice[correct]{\enquote{Bolivianos, el hado propicio}}
		\choice{\enquote{Gloria a Bolivia, tierra bendita}}
		\choice{\enquote{Surge ya la nueva aurora}}
	\end{question}

	
	\begin{question}{Además del coro, el Himno Nacional de Bolivia está compuesto por varias estrofas que desarrollan su mensaje patriótico. ?`Cuántas estrofas tiene oficialmente esta obra musical?}
		\choice{Tres estrofas}
		\choice{Cuatro estrofas}
		\choice{Cinco estrofas}
		\choice{Seis estrofas}
		\choice[correct]{Nueve estrofas}
	\end{question}
\renderSection
